% Ngân hàng tự sinh (AI)
\begin{bt}%[0H9N1-1]
% nguon: AI soạn tương tự câu mẫu 959cabbe
% mau: 959cabbe
	Trên trục tọa độ $\left(O;\overrightarrow{i}\right)$ cho 2 điểm $A$, $B$ có tọa độ lần lượt là $-3$; $4$. Tìm tọa độ của véc-tơ $\overrightarrow{AB}$.
	\loigiai{
		Ta có $\overrightarrow{AB}=4+3=7$.\\
		Vậy toạ độ $\overrightarrow{AB}$ là $7$.
	}
\end{bt}

\begin{bt}%[0H9H1-1]
% nguon: AI soạn tương tự câu mẫu 959cabbe
% mau: 959cabbe nang_cao
	Trên trục tọa độ $\left(O;\overrightarrow{i}\right)$ cho ba điểm $A$, $B$, $C$ có tọa độ lần lượt là $-5$; $2$; $7$. Tìm tọa độ của véc-tơ $\overrightarrow{AC}$ và tọa độ của véc-tơ $\overrightarrow{BC}$.
	\loigiai{
		Ta có $\overrightarrow{AC}=7+5=12$. Vậy toạ độ $\overrightarrow{AC}$ là $12$.\\
		Ta có $\overrightarrow{BC}=7-2=5$. Vậy toạ độ $\overrightarrow{BC}$ là $5$.
	}
\end{bt}

\begin{bt}%[0H9N1-2]
% nguon: AI soạn tương tự câu mẫu ae5e3c24
% mau: ae5e3c24
Trong mặt phẳng tọa độ $Oxy$, cho $2$ véc-tơ $\overrightarrow{a}=(3;4)$, $\overrightarrow{b}=(-8;6)$. Tính tích vô hướng $\overrightarrow{a}\cdot \overrightarrow{b}$ và chứng tỏ $\overrightarrow{a}\perp\overrightarrow{b}$.
\loigiai{
		Ta có $\overrightarrow{a}\cdot \overrightarrow{b} = 3\cdot (-8) + 4\cdot 6 =0$.\\
		Suy ra $\overrightarrow{a}\perp\overrightarrow{b}$.
	}
\end{bt}

\begin{bt}%[0H9H1-2]
% nguon: AI soạn tương tự câu mẫu ae5e3c24
% mau: ae5e3c24 nang_cao
	Trong mặt phẳng tọa độ $Oxy$, cho hai véc-tơ $\overrightarrow{a}=(5;-3)$ và $\overrightarrow{b}=(6;10)$. Tính tích vô hướng $\overrightarrow{a}\cdot \overrightarrow{b}$ và chứng tỏ $\overrightarrow{a}\perp\overrightarrow{b}$.
	\loigiai{
		Ta có $\overrightarrow{a}\cdot \overrightarrow{b} = 5\cdot 6 + (-3)\cdot 10 = 30-30=0$.\\
		Suy ra $\overrightarrow{a}\perp\overrightarrow{b}$.
	}
\end{bt}

\begin{bt}%[0H9N1-3]
% nguon: AI soan tuong tu cau mau 04017960
% mau: 04017960
Trong mặt phẳng tọa độ $Oxy$, cho hai điểm $A(2; -3)$ và $B(5; 4)$. Tìm tọa độ $\overrightarrow{AB}$.
\loigiai{
Ta có $\overrightarrow{AB}=(5-2; 4-(-3))=(3; 7)$.
}
\end{bt}

\begin{bt}%[0H9H1-3]
% nguon: AI soạn tương tự câu mẫu 04017960
% mau: 04017960 nang_cao
	Trong mặt phẳng tọa độ $Oxy$, cho ba điểm $A(-3; 2)$, $B(4; -1)$ và $C(1; 5)$. Gọi $M$ là trung điểm của đoạn thẳng $AB$. Tìm tọa độ $\overrightarrow{MC}$.
	\loigiai{
		Tọa độ trung điểm $M$ của đoạn thẳng $AB$ là $M\left(\dfrac{-3+4}{2}; \dfrac{2+(-1)}{2}\right)$, suy ra $M\left(\dfrac{1}{2}; \dfrac{1}{2}\right)$.\\
		Ta có $\overrightarrow{MC}=\left(1-\dfrac{1}{2}; 5-\dfrac{1}{2}\right)=\left(\dfrac{1}{2}; \dfrac{9}{2}\right)$.\\
		Vậy $\overrightarrow{MC}=\left(\dfrac{1}{2}; \dfrac{9}{2}\right)$.
	}
\end{bt}

\begin{bt}%[0H9H1-4]
% nguon: AI soạn tương tự câu mẫu 2f967af4
% mau: 2f967af4
Trong mặt phẳng $Oxy$, cho $A\left(-2; 1\right), B\left(1; 4\right), C\left(3; -2\right)$. Gọi $M$ là điểm thuộc $BC$ sao cho $\overrightarrow{MB}=-2 \overrightarrow{BC}$ và $G$ là trọng tâm của tam giác $ABC$.
    \begin{enumerate}
        \item Phân tích $\overrightarrow{AM}$ theo hai véc-tơ $\overrightarrow{AB}$ và $\overrightarrow{AC}$. 
        \item Phân tích $\overrightarrow{GM}$ theo hai véc-tơ $\overrightarrow{AB}$ và $\overrightarrow{AC}$.
    \end{enumerate}
	\loigiai{
\begin{enumerate}
        \item Gọi $M(x;y) \Rightarrow \overrightarrow{MB}=(1-x;4-y); \overrightarrow{BC}=(2;-6); \overrightarrow{AB}=(3;3); \overrightarrow{AC}=(5;-3)$.\\
        Mà $\overrightarrow{MB}=-2 \overrightarrow{BC} \Leftrightarrow \heva{&1-x=-4\\ &4-y=12} \Leftrightarrow \heva{&x=5\\ &y=-8} \Rightarrow M(5;-8) \Rightarrow \overrightarrow{AM}=(7;-9)$.\\
        Giả sử $\overrightarrow{AM} = m\overrightarrow{AB}+n\overrightarrow{AC} \Leftrightarrow \heva{&7=3m+5n\\ &-9=3m-3n} \Leftrightarrow \heva{&m= -1\\&n=2.}$\\
        Vậy $\overrightarrow{AM} = -\overrightarrow{AB}+2\overrightarrow{AC}$.
        \item Ta có tọa độ trọng tâm $G\left(\dfrac{2}{3};1\right) \Rightarrow \overrightarrow{GM}=\left(\dfrac{13}{3};-9\right)$.\\
        Giả sử $\overrightarrow{GM} = h\overrightarrow{AB}+k\overrightarrow{AC} \Leftrightarrow \heva{&\dfrac{13}{3}=3h+5k\\ &-9=3h-3k} \Leftrightarrow \heva{&h= -\dfrac{4}{3}\\ &k= \dfrac{5}{3}.}$\\
        Vậy $\overrightarrow{GM} = -\dfrac{4}{3}\overrightarrow{AB}+\dfrac{5}{3}\overrightarrow{AC}$.
    \end{enumerate} 
    \textit{\textbf{Chú ý:} Chúng ta hoàn toàn có thể sử dụng cách biến đổi véc-tơ để thực hiện việc biểu diễn $\overrightarrow{AM},\overrightarrow{GM}$ theo hai véc-tơ $\overrightarrow{AB}$ và $\overrightarrow{AC}$.}
	}
\end{bt}

\begin{bt}%[0H9V1-4]
% nguon: AI soạn tương tự câu mẫu 2f967af4
% mau: 2f967af4 nang_cao
	Trong mặt phẳng $Oxy$, cho $A\left(1; -2\right), B\left(-2; 3\right), C\left(4; 1\right)$. Gọi $M$ là điểm thuộc $BC$ sao cho $\overrightarrow{MB}=-2\overrightarrow{MC}$ và $G$ là trọng tâm của tam giác $ABC$.
	\begin{enumerate}
		\item Phân tích $\overrightarrow{AM}$ theo hai véc-tơ $\overrightarrow{AB}$ và $\overrightarrow{AC}$.
		\item Phân tích $\overrightarrow{GM}$ theo hai véc-tơ $\overrightarrow{AB}$ và $\overrightarrow{AC}$.
	\end{enumerate}
	\loigiai{
		\begin{enumerate}
			\item Gọi $M(x;y) \Rightarrow \overrightarrow{MB}=(-2-x;3-y); \overrightarrow{MC}=(4-x;1-y); \overrightarrow{AB}=(-3;5); \overrightarrow{AC}=(3;3)$.\\
			Mà $\overrightarrow{MB}=-2\overrightarrow{MC} \Leftrightarrow \heva{&-2-x=-8+2x\\ &3-y=-2+2y} \Leftrightarrow \heva{&3x=6\\ &3y=5} \Leftrightarrow \heva{&x=2\\ &y=\dfrac{5}{3}} \Rightarrow M\left(2;\dfrac{5}{3}\right) \Rightarrow \overrightarrow{AM}=\left(1;\dfrac{11}{3}\right)$.\\
			Giả sử $\overrightarrow{AM}=m\overrightarrow{AB}+n\overrightarrow{AC} \Leftrightarrow \heva{&1=-3m+3n\\ &\dfrac{11}{3}=5m+3n} \Leftrightarrow \heva{&m=\dfrac{4}{9}\\ &n=\dfrac{7}{9}.}$\\
			Vậy $\overrightarrow{AM}=\dfrac{4}{9}\overrightarrow{AB}+\dfrac{7}{9}\overrightarrow{AC}$.
			\item Ta có tọa độ trọng tâm $G(1;1) \Rightarrow \overrightarrow{GM}=\left(1;\dfrac{2}{3}\right)$.\\
			Giả sử $\overrightarrow{GM}=h\overrightarrow{AB}+k\overrightarrow{AC} \Leftrightarrow \heva{&1=-3h+3k\\ &\dfrac{2}{3}=5h+3k} \Leftrightarrow \heva{&h=-\dfrac{1}{24}\\ &k=\dfrac{7}{24}.}$\\
			Vậy $\overrightarrow{GM}=-\dfrac{1}{24}\overrightarrow{AB}+\dfrac{7}{24}\overrightarrow{AC}$.
		\end{enumerate}
	}
\end{bt}

\begin{bt}%[0H9N1-5]
% nguon: AI soạn tương tự câu mẫu 42212c6c
% mau: 42212c6c
	Véc-tơ $\overrightarrow{b}=(0;-7)$ được phân tích theo hai véc-tơ đơn vị $\left(\overrightarrow{i};\overrightarrow{j}\right)$ như thế nào?
	\loigiai{
		Ta có $\overrightarrow{b}=(0;-7)\Rightarrow\overrightarrow{b}=0\overrightarrow{i}-7\overrightarrow{j}=-7\overrightarrow{j}$.
	}
\end{bt}

\begin{bt}%[0H9H1-5]
% nguon: AI soạn tương tự câu mẫu 42212c6c
% mau: 42212c6c nang_cao
	Véc-tơ $\overrightarrow{a}=(0;-9)$ được phân tích theo hai véc-tơ đơn vị $\left(\overrightarrow{i};\overrightarrow{j}\right)$ như thế nào? Từ đó tính độ dài của véc-tơ $\overrightarrow{a}$.
	\loigiai{
		Ta có $\overrightarrow{a}=(0;-9)\Rightarrow\overrightarrow{a}=0\overrightarrow{i}-9\overrightarrow{j}=-9\overrightarrow{j}$.\\
		Độ dài véc-tơ: $\left|\overrightarrow{a}\right|=\sqrt{0^2+(-9)^2}=\sqrt{81}=9$.
	}
\end{bt}

\begin{bt}%[0H9H1-6]
% nguon: AI soạn tương tự câu mẫu b279dbb9
% mau: b279dbb9
	Cho $\Delta ABC$ có $A(2; 8)$, $B(6; 1)$, $C(-2; 5)$. Tọa độ điểm $M(a;b)$ thuộc đoạn $BC$ sao cho $S_{\Delta MAB}=3S_{\Delta MAC}$. Tính giá trị của biểu thức $a+b$.
	\loigiai{
		Kẻ $AH \perp BC$. $S_{\Delta MAB}=3S_{\Delta MAC}\Leftrightarrow \dfrac{1}{2}AH\cdot MB=3\cdot\dfrac{1}{2}AH\cdot MC\Leftrightarrow MB=3MC$.\\
		Mà $\overrightarrow{MB}$ và $\overrightarrow{MC}$ ngược hướng $\Rightarrow \overrightarrow{MB}=-3\overrightarrow{MC}$\\
		$\Leftrightarrow \heva{&x_B-x_M=-3(x_C-x_M)\\&y_B-y_M=-3(y_C-y_M)}\Leftrightarrow \heva{&6-x_M=-3(-2-x_M)\\&1-y_M=-3(5-y_M)}$\\
		$\Leftrightarrow \heva{&6-x_M=6+3x_M\\&1-y_M=-15+3y_M}\Leftrightarrow \heva{&4x_M=0\\&4y_M=16}\Rightarrow M(0;4)$.\\
		Suy ra $a+b=4$.
	}
\end{bt}

\begin{bt}%[0H9V1-6]
% nguon: AI soạn tương tự câu mẫu b279dbb9
% mau: b279dbb9 nang_cao
	Cho $\Delta ABC$ có $A(7; 2)$, $B(-3; 5)$, $C(9; -1)$. Tọa độ điểm $M(a;b)$ thuộc đoạn $BC$ sao cho $S_{\Delta MAB}=2S_{\Delta MAC}$. Biết $a+b=\dfrac{m}{n}$ với $\dfrac{m}{n}$ là phân số tối giản. Tính giá trị của biểu thức $T=m+2n$.
	\loigiai{
		Kẻ $AH \perp BC$. $S_{\Delta MAB}=2 S_{\Delta MAC}\Leftrightarrow \dfrac{1}{2}AH\cdot MB=2\cdot \dfrac{1}{2}AH\cdot MC\Leftrightarrow MB=2MC$.\\
		Mà $\overrightarrow{MB}$ và $\overrightarrow{MC}$ ngược hướng $\Rightarrow \overrightarrow{MB}=-2\overrightarrow{MC}$\\
		$\Leftrightarrow \heva{&x_B-x_M=-2(x_C-x_M)\\&y_B-y_M=-2(y_C-y_M)}\Leftrightarrow\heva{&-3-x_M=-2(9-x_M)\\&5-y_M=-2(-1-y_M)}$\\
		$\Leftrightarrow\heva{&-3-x_M=-18+2x_M\\&5-y_M=2+2y_M}\Leftrightarrow\heva{&3x_M=15\\&3y_M=3}\Leftrightarrow\heva{&x_M=5\\&y_M=1}$.\\
		Suy ra $M(5;1)$, do đó $a+b=6=\dfrac{6}{1}$ nên $m=6$, $n=1$.\\
		Vậy $T=m+2n=6+2=8$.
	}
\end{bt}

\begin{bt}%[0H9N2-1]
% nguon: AI soạn tương tự câu mẫu 2dc15a56
% mau: 2dc15a56
	Cho hai véc-tơ $\overrightarrow{a}=(2;1),\overrightarrow{b}=(0;4)$.
	\begin{enumerate}
		\item Tìm toạ độ của véc-tơ $3\overrightarrow{a}+2\overrightarrow{b}$.
		\item Tính các tích vô hướng: $\overrightarrow{a} \cdot \overrightarrow{b},(2\overrightarrow{a}) \cdot (3\overrightarrow{b})$.
	\end{enumerate}
	\loigiai{
		\begin{enumerate}
			\item Ta có: $3\overrightarrow{a}+2\overrightarrow{b}=(3 \cdot 2+2 \cdot 0;3 \cdot 1+2 \cdot 4)=(6;11)$.\\
			\item Ta có: $\overrightarrow{a} \cdot \overrightarrow{b}=2 \cdot 0+1 \cdot 4=4$\\
			$2\overrightarrow{a}=(4;2)$ và $3\overrightarrow{b}=(0;12)$ nên $(2\overrightarrow{a}) \cdot (3\overrightarrow{b})=4 \cdot 0+2 \cdot 12=24$.
		\end{enumerate}
	}
\end{bt}

\begin{bt}%[0H9H2-1]
% nguon: AI soạn tương tự câu mẫu 2dc15a56
% mau: 2dc15a56 nang_cao
	Cho hai véc-tơ $\overrightarrow{a}=(2;-1),\overrightarrow{b}=(1;3)$.
	\begin{enumerate}
		\item Tìm toạ độ của véc-tơ $3\overrightarrow{a}+2\overrightarrow{b}$.
		\item Tính các tích vô hướng: $\overrightarrow{a} \cdot \overrightarrow{b},(2\overrightarrow{a}) \cdot (3\overrightarrow{b})$.
	\end{enumerate}
	\loigiai{
		\begin{enumerate}
			\item Ta có: $3\overrightarrow{a}+2\overrightarrow{b}=(3 \cdot 2+2 \cdot 1;3 \cdot (-1)+2 \cdot 3)=(8;3)$.\\
			\item Ta có: $\overrightarrow{a} \cdot \overrightarrow{b}=2 \cdot 1+(-1) \cdot 3=-1$\\
			$2\overrightarrow{a}=(4;-2)$ và $3\overrightarrow{b}=(3;9)$ nên $(2\overrightarrow{a}) \cdot (3\overrightarrow{b})=4 \cdot 3+(-2) \cdot 9=-6$.
		\end{enumerate}
	}
\end{bt}

\begin{bt}%[0H9N2-2]
% nguon: AI soạn tương tự câu mẫu 34f6d564
% mau: 34f6d564
	Cho $\overrightarrow{a}=(3; 2)$ và $\overrightarrow{b}=(-2;5)$. Tính $\cos \alpha$, với $\alpha =\left(\overrightarrow{a}, \overrightarrow{b}\right)$.
	\loigiai{
		Ta có $\cos \alpha =\dfrac{\overrightarrow{a}\cdot \overrightarrow{b}}{\left|\overrightarrow{a}\right|\cdot \left|\overrightarrow{b}\right|} =\dfrac{3\cdot (-2)+2\cdot 5}{\sqrt{3^2+2^2}\cdot \sqrt{(-2)^2+5^2}}=\dfrac{4}{\sqrt{377}}$.
	}
\end{bt}

\begin{bt}%[0H9H2-2]
% nguon: AI soạn tương tự câu mẫu 34f6d564
% mau: 34f6d564 nang_cao
	Cho $\overrightarrow{a}=(3; -1)$ và $\overrightarrow{b}=(2;5)$. Tính $\cos \alpha$, với $\alpha =\left(\overrightarrow{a}, \overrightarrow{b}\right)$.
	\loigiai{
		Ta có $\cos \alpha =\dfrac{\overrightarrow{a}\cdot \overrightarrow{b}}{\left|\overrightarrow{a}\right|\cdot \left|\overrightarrow{b}\right|} =\dfrac{3\cdot 2+(-1)\cdot 5}{\sqrt{3^2+(-1)^2}\cdot \sqrt{2^2+5^2}}=\dfrac{1}{\sqrt{290}}$.\\
		Vậy $\cos \alpha =\dfrac{1}{\sqrt{290}}$.
	}
\end{bt}

\begin{bt}%[0H9H2-3]
% nguon: AI soạn tương tự câu mẫu a3548811
% mau: a3548811
	Trong mặt phẳng tọa độ $Oxy$, cho tam giác $ABC$ có $A(2;-1)$, $B(3;3)$ và $C(9;1)$. Tính chu vi tam giác $ABC$.
	\loigiai{
		Ta có $\heva{&\overrightarrow{AB}=(1;4)\\&\overrightarrow{AC}=(7;2)\\&\overrightarrow{BC}=(6;-2)} \Rightarrow \heva{&AB=\sqrt{1^2+4^2}=\sqrt{17}\\&AC=\sqrt{7^2+2^2}=\sqrt{53}\\&BC=\sqrt{6^2+(-2)^2}=2\sqrt{10}.}$\\
		Chu vi tam giác $ABC$ là $AB+AC+BC=\sqrt{17}+\sqrt{53}+2\sqrt{10}$.
	}
\end{bt}

\begin{bt}%[0H9V2-3]
% nguon: AI soạn tương tự câu mẫu a3548811
% mau: a3548811 nang_cao
	Trong mặt phẳng tọa độ $Oxy$, cho tam giác $ABC$ có $A(2;-3)$, $B(5;1)$ và $C(11;-1)$. Gọi $M$ là trung điểm của cạnh $AC$. Tính chu vi tam giác $ABM$.
	\loigiai{
		Vì $M$ là trung điểm của $AC$ nên $M\left(\dfrac{2+11}{2};\dfrac{-3+(-1)}{2}\right)=\left(\dfrac{13}{2};-2\right)$.\\
		Ta có $\heva{&\overrightarrow{AB}=(3;4)\\&\overrightarrow{AM}=\left(\dfrac{9}{2};1\right)\\&\overrightarrow{BM}=\left(\dfrac{3}{2};-3\right)} \Rightarrow \heva{&AB=\sqrt{3^2+4^2}=5\\&AM=\sqrt{\left(\dfrac{9}{2}\right)^2+1^2}=\dfrac{\sqrt{85}}{2}\\&BM=\sqrt{\left(\dfrac{3}{2}\right)^2+(-3)^2}=\dfrac{3\sqrt{5}}{2}.}$\\
		Chu vi tam giác $ABM$ là $AB+AM+BM=5+\dfrac{\sqrt{85}}{2}+\dfrac{3\sqrt{5}}{2}=\dfrac{10+\sqrt{85}+3\sqrt{5}}{2}$.
	}
\end{bt}

\begin{bt}%[0H9N2-4]
% nguon: AI soạn tương tự câu mẫu 0e8eb5bf
% mau: 0e8eb5bf
	Trong mặt phẳng tọa độ $Oxy$, cho điểm $A(2;-3)$, đường thẳng $\Delta_1\colon 2x-y+5=0$ và đường thẳng $\Delta_2\colon x+3y-10=0$.
	\begin{enumerate}
		\item Tính góc giữa đường thẳng $\Delta_1$ và $\Delta_2$.
		\item Tính khoảng cách từ điểm $A$ đến đường thẳng $\Delta_2$.
	\end{enumerate}
	\loigiai{
		\begin{enumerate}
			\item $\Delta_1$ có véc-tơ pháp tuyến $\overrightarrow{n}_1=(2;-1)$ và $\Delta_2$ có véc-tơ pháp tuyến $\overrightarrow{n}_2=(1;3)$.\\
			$\cos(\Delta_1;\Delta_2)=\left|\cos(\overrightarrow{n}_1;\overrightarrow{n}_2)\right|=\dfrac{|2\cdot1+(-1)\cdot3|}{\sqrt{2^2+(-1)^2}\cdot\sqrt{1^2+3^2}}=\dfrac{1}{\sqrt{5}\cdot\sqrt{10}}=\dfrac{\sqrt{2}}{10}$.\\
			Suy ra $(\Delta_1;\Delta_2)\approx 81{,}87^\circ$.
			\item $\mathrm{d}(A,\Delta_2)=\dfrac{|1\cdot2+3\cdot(-3)-10|}{\sqrt{1^2+3^2}}=\dfrac{17}{\sqrt{10}}=\dfrac{17\sqrt{10}}{10}$.
		\end{enumerate}
	}
\end{bt}

\begin{bt}%[0H9H2-4]
% nguon: AI soạn tương tự câu mẫu 0e8eb5bf
% mau: 0e8eb5bf nang_cao
Trong mặt phẳng tọa độ $Oxy$, cho điểm $A(1;-2)$, đường thẳng $\Delta_1\colon 3x+y-7=0$ và đường thẳng $\Delta_2\colon x-2y+4=0$.
	\begin{enumerate}
		\item Tính góc giữa đường thẳng $\Delta_1$ và $\Delta_2$.
		\item Tính khoảng cách từ điểm $A$ đến đường thẳng $\Delta_2$.
	\end{enumerate}
	\loigiai{
		\begin{enumerate}
			\item $\Delta_1$ có véc-tơ pháp tuyến $\overrightarrow{n}_1=(3;1)$ và $\Delta_2$ có véc-tơ pháp tuyến $\overrightarrow{n}_2=(1;-2)$.\\
			$\cos(\Delta_1;\Delta_2)=\left|\cos(\overrightarrow{n}_1;\overrightarrow{n}_2)\right|=\dfrac{|3\cdot1+1\cdot(-2)|}{\sqrt{3^2+1^2}\cdot\sqrt{1^2+(-2)^2}}=\dfrac{1}{\sqrt{10}\cdot\sqrt{5}}=\dfrac{1}{5\sqrt{2}}=\dfrac{\sqrt{2}}{10}$.\\
			Suy ra $(\Delta_1;\Delta_2)\approx 81{,}87^\circ$.
			\item $\mathrm{d}(A,\Delta_2)=\dfrac{|1\cdot1-2\cdot(-2)+4|}{\sqrt{1^2+(-2)^2}}=\dfrac{9}{\sqrt{5}}=\dfrac{9\sqrt{5}}{5}$.
		\end{enumerate}
	}
\end{bt}

\begin{bt}%[0H9N2-5]
% nguon: AI soạn tương tự câu mẫu 4924d658
% mau: 4924d658
Trong mặt phẳng tọa độ $Oxy$, cho $\triangle ABC$ có $A(0;6), B(-6;3), C(2;3)$.
\begin{enumerate}
	\item Xác định tọa độ trung điểm $I$ của $AB$ và trọng tâm $G$ của $\triangle ABC$.
	\item Viết phương trình tổng quát của đường thẳng $AB$.
	\item Viết phương trình đường tròn $(C)$ có tâm $A$ và đi qua $C$.
	\item Viết phương trình đường thẳng $d$ qua điểm $C$ đồng thời cắt $Ox$ và $Oy$ lần lượt tại hai điểm $E$ và $F$ sao cho $\triangle OEF$ có diện tích bằng $1$.
\end{enumerate}
	\loigiai{
\begin{enumerate}
	\item Trung điểm của $AB$ là $I(-3;4{,}5)$.\\
	Trọng tâm của $\triangle ABC$ là $G\left(-\dfrac{4}{3};4\right)$.
	\item Đường thẳng $AB$ có véc-tơ chỉ phương là $\vec{u}=\vec{AB}=(-6;-3)$.\\
	$AB$ đi qua $A$ và có véc-tơ pháp tuyến $\vec{n}=(1;-2)$ nên $AB$ có phương trình là
	$$1(x-0)-2(y-6)=0 \Leftrightarrow x-2y+12=0.$$
	\item Vì $(C)$ có tâm $A$ và đi qua $C$ nên $(C)$ có bán kính là $R=AC=\sqrt{(2-0)^2+(3-6)^2}=\sqrt{13}$.\\
	Vậy $(C)$ có phương trình là $x^2+(y-6)^2=13$.
	\item Gọi $E(a;0) \in Ox$ và $F(0;b)\in Oy$.\\
	Vì $d$ qua $E$ và $F$ nên phương trình của $d$ theo đoạn chắn là $\dfrac{x}{a}+\dfrac{y}{b}=1$.\\
	Mà $d$ qua $C(2;3)$ nên ta có $\dfrac{2}{a}+\dfrac{3}{b}=1 \Leftrightarrow 3a+2b=a\cdot b$ \qquad (1).\\
	$S_{\triangle OEF}= 1 \Leftrightarrow \dfrac{1}{2} OE \cdot OF = 1 \Leftrightarrow \dfrac{1}{2} |a|\cdot |b| =1 \Leftrightarrow |a\cdot b| =2 \Leftrightarrow \hoac{&a\cdot b= 2 \,\,\quad (2) \\& a\cdot b= -2 \quad (3).}$\\
	Từ (1) và (2) ta có $\heva{&3a+2b=2 \\&a\cdot b=2} \Leftrightarrow \heva{&b=\dfrac{2-3a}{2} \\& a\cdot\dfrac{2-3a}{2}=2} \Leftrightarrow \heva{&b=\dfrac{2-3a}{2} \\&-3a^2+2a-4=0.}$ (hệ vô nghiệm).\\
	Từ (1) và (3) ta có $\heva{&3a+2b=-2 \\&a\cdot b=-2} \Leftrightarrow \heva{&b=\dfrac{-2-3a}{2} \\& a\cdot\dfrac{-2-3a}{2}=-2} \Leftrightarrow \heva{&b=\dfrac{-2-3a}{2} \\&3a^2+2a-4=0} \Leftrightarrow \heva{& b=\dfrac{-2-3a}{2} \\& \hoac{&a=\dfrac{-1+\sqrt{13}}{3} \\& a= \dfrac{-1-\sqrt{13}}{3}.}}$\\
	$\bullet$ Với $a=\dfrac{-1+\sqrt{13}}{3} \Rightarrow b=\dfrac{-1-\sqrt{13}}{2}$ khi đó $d\colon \dfrac{3x}{-1+\sqrt{13}}+\dfrac{2y}{-1-\sqrt{13}}=1$.\\
	$\bullet$ Với $a=\dfrac{-1-\sqrt{13}}{3} \Rightarrow b=\dfrac{-1+\sqrt{13}}{2}$ khi đó $d\colon \dfrac{3x}{-1-\sqrt{13}}+\dfrac{2y}{-1+\sqrt{13}}=1$.\\
	Vậy có hai đường thẳng $d$ thỏa đề bài.
	% KIEM_TRA: Đề mẫu có đáp số đẹp vì chọn C(1;2) và S=1/2; với C(2;3) và S=1 hệ số a,b chứa căn, không "đẹp" như mẫu nhưng vẫn đúng cấu trúc và duy nhất.
\end{enumerate}
	}
\end{bt}

\begin{bt}%[0H9H2-5]
% nguon: AI soạn tương tự câu mẫu 4924d658
% mau: 4924d658 nang_cao
Trong mặt phẳng tọa độ $Oxy$, cho $\triangle ABC$ có $A(0;6)$, $B(-6;3)$, $C(2;3)$.
\begin{enumerate}
	\item Xác định tọa độ trung điểm $I$ của $AB$ và trọng tâm $G$ của $\triangle ABC$.
	\item Viết phương trình tổng quát của đường thẳng $AB$.
	\item Viết phương trình đường tròn $(C)$ có tâm $A$ và đi qua $C$.
	\item Viết phương trình đường thẳng $d$ qua điểm $C$ đồng thời cắt $Ox$ và $Oy$ lần lượt tại hai điểm $E$ và $F$ sao cho $\triangle OEF$ có diện tích bằng $1$.
\end{enumerate}
	\loigiai{
		\begin{enumerate}
			\item Trung điểm của $AB$ là $I(-3;4{,}5)$.\\
			Trọng tâm của $\triangle ABC$ là $G\left(-\dfrac{4}{3};4\right)$.
			\item Đường thẳng $AB$ có véc-tơ chỉ phương là $\vec{u}=\vec{AB}=(-6;-3)$.\\
			$AB$ đi qua $A$ và có véc-tơ pháp tuyến $\vec{n}=(1;-2)$ nên $AB$ có phương trình là
			$$1(x-0)-2(y-6)=0 \Leftrightarrow x-2y+12=0.$$
			\item Vì $(C)$ có tâm $A$ và đi qua $C$ nên $(C)$ có bán kính là $R=AC=\sqrt{(2-0)^2+(3-6)^2}=\sqrt{13}$.\\
			Vậy $(C)$ có phương trình là $x^2+(y-6)^2=13$.
			\item Gọi $E(a;0) \in Ox$ và $F(0;b)\in Oy$.\\
			Vì $d$ qua $E$ và $F$ nên phương trình của $d$ theo đoạn chắn là $\dfrac{x}{a}+\dfrac{y}{b}=1$.\\
			Mà $d$ qua $C(2;3)$ nên ta có $\dfrac{2}{a}+\dfrac{3}{b}=1 \Leftrightarrow 3a+2b=a\cdot b$ \qquad (1).\\
			$S_{\triangle OEF}= 1 \Leftrightarrow \dfrac{1}{2} OE \cdot OF = 1 \Leftrightarrow \dfrac{1}{2} |a|\cdot |b| =1 \Leftrightarrow |a\cdot b| =2 \Leftrightarrow \hoac{&a\cdot b= 2 \,\,\quad (2) \\& a\cdot b= -2 \quad (3).}$\\
			Từ (1) và (2) ta có $\heva{&3a+2b=2 \\&a\cdot b=2} \Leftrightarrow \heva{&b=\dfrac{2-3a}{2} \\& a\cdot\dfrac{2-3a}{2}=2} \Leftrightarrow \heva{&b=\dfrac{2-3a}{2} \\&3a^2-2a+4=0}$ (vô nghiệm vì $\Delta'=1-12<0$).\\
			Từ (1) và (3) ta có $\heva{&3a+2b=-2 \\&a\cdot b=-2} \Leftrightarrow \heva{&b=\dfrac{-2-3a}{2} \\& a\cdot\dfrac{-2-3a}{2}=-2} \Leftrightarrow \heva{&b=\dfrac{-2-3a}{2} \\&3a^2+2a-4=0} \Leftrightarrow \hoac{&a=\dfrac{-1+\sqrt{13}}{3} \\& a=\dfrac{-1-\sqrt{13}}{3}.}$\\
			$\bullet$ Với $a=\dfrac{-1+\sqrt{13}}{3}$ thì $b=\dfrac{-3-\sqrt{13}}{2}$, khi đó $d\colon \dfrac{3x}{-1+\sqrt{13}}+\dfrac{2y}{-3-\sqrt{13}}=1$.\\
			$\bullet$ Với $a=\dfrac{-1-\sqrt{13}}{3}$ thì $b=\dfrac{-3+\sqrt{13}}{2}$, khi đó $d\colon \dfrac{3x}{-1-\sqrt{13}}+\dfrac{2y}{-3+\sqrt{13}}=1$.\\
			Vậy có hai đường thẳng $d$ thỏa đề bài.
		\end{enumerate}
	}
\end{bt}
